\section{Actual directional derivatives and first dehomogeneous normal bands}

Fix the actual point chart $(e,b)$, write $S$ for its genuine source
linear shear, and put $v_e=1$, $v_i=b_i$ for $i\ne e$.
The derivation $\delta_v=\sum_i v_i\partial_i$ is literal polynomial
differentiation. Normality always means coefficient vanishing on the
ENTIRE original source diagram, not on a newly chosen diagram.

\begin{theorem}[Directional differentiation preserves source normality and computes actual chart bands]
\label{hir:actual-directional-normal-polynomial-dehom-euler}
\uses{hir:actual-formal-ideal-initial-diagram-upper,hir:actual-point-exceptional-restriction-dehomogenized-polynomial}
\lean{TranscendentalNakamaye.Pending.HironakaActualDirectionalNormalJet.directionalPolynomialDerivative,TranscendentalNakamaye.Pending.HironakaActualDirectionalNormalJet.directional_derivative_normal,TranscendentalNakamaye.Pending.HironakaActualDirectionalNormalJet.directional_derivative_homogeneous,TranscendentalNakamaye.Pending.HironakaActualDirectionalNormalJet.dehom_directional_derivative_coefficient}
\leanok
The literal directional derivative of a source-diagram-normal polynomial
is again source-diagram normal. If $P$ is homogeneous of degree $d$ and
$f=\operatorname{dehom}_{e,b}P$, then at every exponent $\beta$,
\[
 [X^\beta]\operatorname{dehom}_{e,b}(\delta_vP)
   =(d-|\beta|)[X^\beta]f.
\]
The scalar difference is taken in $\mathbb C$. For $d>0$ the derivative is homogeneous of degree $d-1$;
for $d=0$ it is zero.
\end{theorem}
\begin{proof}
At a source-diagram exponent $\beta$, each derivative coefficient is
$(\beta_i+1)P_{\beta+\mathbf1_i}$. Upward closure puts every latter
exponent in the source diagram, so each summand is zero. This proves
normality for the finite actual directional sum.

For $i\ne e$, direct differentiation of each actual substitution
$X_e\mapsto1$, $X_j\mapsto b_j+X_j$ shows
$\partial_i f=\operatorname{dehom}(\partial_iP)$.
The same identity follows for products by the Leibniz rule and for sums
and constants by additivity; thus it holds for the entire polynomial.
Apply dehomogenization to Euler's homogeneous identity
$\sum_i X_i\partial_iP=dP$. Split the pivot term from the other
actual coordinates and collect the constant factors $b_i$. It gives
\[
 \operatorname{dehom}(\delta_vP)+\sum_i X_i\partial_i f=df.
\]
For each monomial of $f$, the second sum multiplies its coefficient by
the sum of its exponents, namely $|\beta|$. Taking its genuine
coefficient gives the displayed identity. Each partial derivative of a
degree-$d$ homogeneous polynomial is homogeneous of degree $d-1$,
including the zero case, and finite constant linear combinations retain
that homogeneity.
\end{proof}

\begin{theorem}[A first actual dehomogeneous normal band cannot belong to the sheared source initial ideal]
\label{hir:actual-first-dehom-normal-band-initial-independence}
\uses{hir:actual-directional-normal-polynomial-dehom-euler,hir:actual-point-homogeneous-dehom-shear-identity,hir:actual-formal-normal-representative-unique,hir:actual-homogeneous-initial-ideal-diagram-equality}
\lean{TranscendentalNakamaye.Pending.HironakaActualDirectionalNormalJet.first_dehomogeneous_normal_band_eq_zero}
\leanok
Let $K$ be an actual whole source formal ideal and $P$ a homogeneous
degree-$d$ polynomial normal on $\Gamma(K)$. Let $k\le d$, and assume
the actual dehomogenization $f$ has every coefficient of degree below
$k$ zero. If its whole degree-$k$ band belongs to the genuine sheared
ideal $S(\operatorname{in}K)$, that band is zero.
\end{theorem}
\begin{proof}
Induct on $d-k$. When $d=k$, the actual affine-linear dehomogenization
has degree at most $d$, and all lower coefficients vanish. It is
therefore genuinely homogeneous of degree $d$. The actual homogeneous
dehomogenization/shear identity gives $S(P)=f$. The membership assumption
and the genuine inverse shear imply $P\in\operatorname{in}K$.
The whole initial diagram equals $\Gamma(K)$, whereas all coefficients
of $P$ there vanish. Actual normal-representative uniqueness makes
$P=0$, and hence its actual dehomogeneous band is zero.

When $d>k$, apply the preceding directional derivative theorem. The
polynomial $\delta_vP$ is homogeneous of degree $d-1$ and source
normal. Its actual dehomogenization still has all bands below $k$ zero,
and its degree-$k$ band is exactly $(d-k)f_k$. This band belongs to the
same whole sheared ideal because $f_k$ does. The induction hypothesis
makes it zero. The positive integer $d-k$ is a nonzero complex scalar,
so scalar cancellation yields $f_k=0$. This argument retains the
original diagram and uses no source/current diagram inclusion.
\end{proof}

\begin{theorem}[The same canonical higher source bands supply the normal input]
\label{hir:actual-canonical-higher-band-dehom-initial-independence}
\uses{hir:actual-first-dehom-normal-band-initial-independence,hir:actual-reduced-formal-vertex-members}
\lean{TranscendentalNakamaye.Pending.HironakaActualDirectionalNormalJet.canonical_higher_dehomogeneous_band_eq_zero}
\leanok
For every SAME canonical reduced source series $G_a$, every degree
$d>|a|$ source band satisfies the preceding independence statement
internally, without an incoming source-normality premise.
\end{theorem}
\begin{proof}
Take the actual finite degree-$d$ homogeneous polynomial of $G_a$.
At a source-diagram exponent its coefficient equals the corresponding
whole canonical reduced coefficient. Such a coefficient is zero except
possibly at $a$. That exception cannot have degree $d>|a|$.
Thus the entire actual source band is source normal. It is genuinely
homogeneous by construction. The preceding theorem applies to its
actual dehomogeneous bands and the explicitly stated whole sheared
initial-ideal membership, proving the conclusion for this SAME source
family.
\end{proof}

The chart-band vanishing and ideal membership here are explicit finite
algebra hypotheses. Their derivation from a first offending current
normalized band and original persistent point geometry is separate.
No full arbitrary-direction or positive-centre HS criterion is asserted.
